%% ============================================================
%% CHAPTER 3 — METHODOLOGY
%% ============================================================
\chapter{Methodology}
\label{ch:method}

\begin{figure}[H]
\centering
\begin{tikzpicture}[
    node distance=0.8cm,
    box/.style={rectangle, draw=uqpurple, fill=uqpurple!5, thick, rounded corners, text width=0.9\textwidth, minimum height=1.2cm, align=center, font=\small, inner sep=1ex},
    arrow/.style={-{Stealth[scale=1.2]}, thick, uqpurple}
]

\node[box] (lit) {\textbf{Phase 1: Knowledge Gap \& Validation} \\ Literature review, parameter selection, and CoolProp/NIST model validation.};

\node[box, below=of lit] (thermo) {\textbf{Phase 2: Thermodynamic \& Exergy Modelling (RQ1)} \\ First-law and second-law nodal analysis. Multi-point T-H pinch extraction.};

\node[box, below=of thermo] (sizing) {\textbf{Phase 3: Component Sizing \& Capital Cost Estimation (RQ2)} \\ NTU-effectiveness heat exchanger sizing (with pseudo-critical correction) and turbine VFR/SP calculation.};

\node[box, below=of sizing] (econ) {\textbf{Phase 4: Techno-Economic \& Off-Design Analysis (RQ3 \& RQ4)} \\ LCOE/NPV modelling, 10,000-sample Monte Carlo uncertainty, and Stodola part-load degradation.};

\draw[arrow] (lit) -- (thermo);
\draw[arrow] (thermo) -- (sizing);
\draw[arrow] (sizing) -- (econ);

\end{tikzpicture}
\caption{Overview of the interconnected four-phase methodology used to address RQ1--RQ4.}
\label{fig:methodology_flow}
\end{figure}

The research was structured in four phases corresponding to the four research
questions. The approach was set out in the thesis proposal, substantially
expanded in the Interim Report, and revised during the final modelling work
described here. Three revisions made after the Interim Report are material to
the results and are flagged where they arise: (i)~the recuperator energy
balance is limited by both streams (Section~\ref{sec:method_recup}), which
lowers the CO$_2$ cycle efficiencies reported at the interim stage;
(ii)~every primary heater is integrated with the waste-gas stream using a
duty-linear gas profile and a 10\,K approach (Section~\ref{sec:method_pinch});
and (iii)~the techno-economic comparison is made on a fixed-source basis
(Section~\ref{sec:method_rq3}).

\section{Phase 1 --- Literature Review}
\label{sec:method_lit}

A structured search was conducted across Scopus, Web of Science, and Google
Scholar following the PRISMA framework, covering WHR publications from 2005
to 2025 with five foundational pre-2005 sources included on merit (Dostal
2004; Short et al.\ 1995; Moran \& Shapiro 2018; Turton et al.\ 2012; Bell
et al.\ 2014). Coverage was organised around six themes: cycle thermodynamics
and variants; experimental and simulation performance data; transitional and
pseudo-critical CO$_2$ thermophysics; heat exchanger technology for
high-pressure CO$_2$; capital and operating cost data; and Australian
industrial waste heat availability. Thirty-two references were retained after
relevance screening; 29 are cited in the body of the thesis.

\section{Thermodynamic Modelling (RQ1)}
\label{sec:method_thermo}

\subsection{Cycle Architecture and State-Point Models}

Steady-state energy balance models for all three cycles were implemented in
Python\,3.9 using the CoolProp\,6.x library~\cite{Bell2014} for
thermodynamic property evaluation. The HEOS (Helmholtz Energy Equation of
State) backend was used throughout in preference to the default BICUBIC
interpolation because of its superior accuracy near the CO$_2$ critical
point, where the tCO$_2$ pump exit and the sCO$_2$ compressor inlet both
operate. At each state point four properties are extracted: specific enthalpy
($h$), specific entropy ($s$), density ($\rho$), and temperature ($T$).

Energy balances at each state point follow standard open-system forms. For
turbomachinery:
\begin{equation}
\dot{W} = \dot{m}(h_{in} - h_{out})
\label{eq:turbomachinery}
\end{equation}
For heat exchangers:
\begin{equation}
\dot{Q} = \dot{m}(h_{out} - h_{in})
\label{eq:hx}
\end{equation}
Isentropic efficiency is applied as:
\begin{equation}
h_{out} = h_{in} + \frac{h_{out,is} - h_{in}}{\eta_{is}}
\quad\text{(compression)}
\label{eq:compression}
\end{equation}
\begin{equation}
h_{out} = h_{in} - \eta_{is}(h_{in} - h_{out,is})
\quad\text{(expansion)}
\label{eq:expansion}
\end{equation}
Values of $\eta_{turbine} = 0.85$ and $\eta_{pump/compressor} = 0.80$ are
applied uniformly to all three cycles, ensuring genuine comparability. These
values are standard for preliminary analysis of small-to-medium scale
turbomachinery~\cite{Macchi2017} and represent a known simplification
addressed further in Section~\ref{sec:discussion_limitations}.

All heat exchangers and pipe runs are modelled as isobaric processes
($\Delta P = 0$) in the thermodynamic analysis phase. Real frictional pressure
drop --- particularly significant in CO$_2$ PCHEs --- is an identified
limitation acknowledged and quantified in Section~\ref{sec:discussion_limitations}.

\subsection{Cycle-Specific Parametric Assumptions}

The tCO$_2$ high-side pressure is scaled linearly from 12\,MPa at
150\,\textdegree C to 20\,MPa at 350\,\textdegree C and above, following
the optimal pressure relationship of Liu et al.~\cite{Liu2024}
($P_{high} = \min[20, \max(12,\, 8 + 0.048\,(T-100))]$\,MPa). Low-side
pressure is held at 5\,MPa for the baseline, corresponding to a condenser
temperature of 14.3\,\textdegree C. The recuperator effectiveness baseline is
$\varepsilon = 0.70$ for tCO$_2$ and $\varepsilon = 0.85$ for sCO$_2$.

\phantomsection\label{sec:method_recup}
The recuperator effectiveness is applied to the stream with the smaller
enthalpy capacity, so that neither stream can be driven past the inlet
temperature of the other:
\begin{equation}
Q_{recup} = \varepsilon \cdot \min\!\big[\,h_{cold}(T_{hot,in}) - h_{cold,in},\;
h_{hot,in} - h_{hot}(T_{cold,in})\,\big]
\label{eq:recup}
\end{equation}
This is the standard $\varepsilon$-NTU definition~\cite{Moran2018}. It
matters for CO$_2$ because $c_p$ differs strongly between the high- and
low-pressure sides near the critical point: at the CS2 condition the cold
(12.8\,MPa) side can absorb 251\,kJ/kg before reaching the turbine exhaust
temperature, whereas the hot (5\,MPa) side can release only 113\,kJ/kg before
reaching the pump exit temperature. Defining $\varepsilon$ on the cold side
alone, as was done in the Interim Report, transfers 176\,kJ/kg and drives the
hot stream to 14\,\textdegree C --- 10\,K \emph{below} the stream it is
heating --- which violates the second law. Equation~\eqref{eq:recup} is used
throughout this thesis, and the resulting efficiencies are lower than the
interim values (Section~\ref{sec:results_firstlaw}). The recuperator is also
discretised in 200 enthalpy steps to confirm that no internal temperature
crossover occurs (Section~\ref{sec:method_rq2}).

The sCO$_2$ cycle uses $P_{high} = 20$\,MPa and $P_{low} = 8$\,MPa, with
the compressor inlet set to $T_c + 5$\,K (35.98\,\textdegree C) to guarantee
operation in the supercritical region and maintain the near-critical
compressibility benefit.

In the parametric sweep the SRC uses $P_{high} = \min[10\,\text{MPa},\;
0.95 \times P_{sat}(T_{TIT})]$, i.e.\ the highest subcritical pressure that
leaves the turbine inlet slightly superheated, and $P_{low} = 10$\,kPa
(condenser temperature 45.8\,\textdegree C), a typical industrial condensing
pressure. This sweep characterises each cycle as a function of turbine inlet
temperature only; when the cycles are integrated with a real waste-gas stream
(Section~\ref{sec:method_rq2}) the steam pressure is instead chosen by the
boiler pinch constraint, which gives substantially lower pressures.

\subsection{Model Validation}

CoolProp outputs were cross-validated against NIST Webbook
data~\cite{Lemmon2018} at two SRC state points (Table~\ref{tab:validation}).
Both deviations lie well within the 1\,\% acceptance criterion. CoolProp
evaluates the CO$_2$ critical point as $T_c = 30.98$\,\textdegree C and
$P_c = 7.377$\,MPa, matching NIST REFPROP to within 0.04\,\%.

\begin{table}[H]
\centering
\caption{SRC model validation against NIST Webbook reference data.}
\label{tab:validation}
\renewcommand{\arraystretch}{1.2}
\begin{tabularx}{\textwidth}{>{\raggedright\arraybackslash}X
  >{\centering\arraybackslash}p{2.5cm}
  >{\centering\arraybackslash}p{2.5cm}
  >{\centering\arraybackslash}p{2.5cm}
  >{\centering\arraybackslash}p{2.0cm}}
\toprule
\rowcolor{tableheader}
\textbf{State point} & \textbf{NIST $h$} & \textbf{Model $h$} &
\textbf{NIST $s$} & \textbf{Deviation} \\
\rowcolor{tableheader}
& \textbf{(kJ/kg)} & \textbf{(kJ/kg)} & \textbf{(kJ/kg$\cdot$K)} & \\
\midrule
Saturated liquid, 10\,kPa & 191.83 & 191.81 & 0.6493 & 0.010\,\% \\
\rowcolor{rowgray}
Superheated steam, 400\,\textdegree C, 10\,MPa
& 3096.5 & 3097.45 & 6.2120 & 0.031\,\% \\
\bottomrule
\end{tabularx}
\end{table}

The global exergy balance (Eq.~\ref{eq:exergy_balance}) closes to machine
precision ($<10^{-12}$\,kJ/kg) for all three cycles, confirming internal
consistency of the second-law model.

\subsection{Parametric Sweep}

The primary sweep covers TIT from 150 to 600\,\textdegree C in
25\,\textdegree C steps. At each temperature, $\eta_{th}$, net specific
work, BWR, and exergy destruction by component are recorded. Separate sweeps
for recuperator effectiveness (Section~\ref{sec:method_sensitivity}) and
low-side pressure (Section~\ref{sec:method_sensitivity}) extend the analysis
beyond the single-parameter baseline.

\section{Exergy Analysis (RQ1)}
\label{sec:method_exergy}

An exergy balance is applied to every cycle component at a common turbine
inlet temperature of 400\,\textdegree C and at the three site design points
(Section~\ref{sec:method_rq2}). The dead state is $T_0 = 298.15$\,K,
$P_0 = 101.325$\,kPa (ASHRAE standard, confirmed with supervisor). Because
the three cycles circulate very different mass flows per unit of output,
component destruction is reported both per kilogram of working fluid and
per kilojoule of net work; only the latter is a like-for-like basis.

Specific flow exergy:
\begin{equation}
ex = (h - h_0) - T_0(s - s_0)
\label{eq:flow_exergy}
\end{equation}
For adiabatic turbomachinery:
\begin{equation}
\dot{E}_{x,dest} = T_0\,\dot{m}(s_{out} - s_{in})
\label{eq:adiabatic_exdest}
\end{equation}
For heat exchangers receiving heat from source at $T_{source}$:
\begin{align}
ex_{heat,in} &= Q_{in}\!\left(1 - \frac{T_0}{T_{source}}\right) \label{eq:exergy_heat}\\
\dot{E}_{x,dest,HX} &= ex_{heat,in} - (ex_{out} - ex_{in}) \label{eq:hx_exdest}
\end{align}
For condensers rejecting heat to environment:
\begin{equation}
\dot{E}_{x,dest,cond} = ex_{in} - ex_{out}
\label{eq:cond_exdest}
\end{equation}
Global balance closure check:
\begin{equation}
\dot{W}_{net} + \dot{E}_{x,dest,total} = ex_{heat,in}
\label{eq:exergy_balance}
\end{equation}

\section{Sensitivity Analysis (RQ1)}
\label{sec:method_sensitivity}

\subsection{Recuperator Effectiveness Sweep}

The tCO$_2$ cycle is evaluated at four recuperator effectiveness values
($\varepsilon = 0, 0.50, 0.70, 0.85$) at each of the three site design
points, with all other parameters held at baseline. For each $\varepsilon$
value the heater integration of Section~\ref{sec:method_pinch} is repeated to
determine how recuperation raises the heater inlet temperature and therefore
the lowest temperature to which the waste gas can be cooled. The quantity of
interest for waste heat recovery is not efficiency alone but the net work
obtained per kilogram of waste gas, $w_{gas} = \eta_{th}\,c_{p,gas}\,
(T_{src} - T_{gas,out})$, which combines both effects.

\subsection{Low-Side Pressure Sweep}

The tCO$_2$ cycle is evaluated at $P_{low} = 3.5, 4.0, 5.0, 6.0$ and
$7.0$\,MPa with $\varepsilon = 0.70$ fixed, at each site design point.
Condensing temperature, $\eta_{th}$, BWR, achievable gas outlet temperature
and $w_{gas}$ are recorded. The feasibility of each $P_{low}$ value is assessed
against the cooling-water temperatures available at the three Australian
sites, and the parasitic load of mechanical chilling, where it would be
required, is estimated from the condenser duty and a chiller coefficient of
performance of 4--6.

